% s-039-p039-m002; original PDF 39, 40
\Aphorism{s-039-p039-m002}{24}
Es ist ein Trost, dass Einsamkeit und Unabhängigkeit zu den unerreichbaren Dingen gehören; man soll sie wohl suchen, würde aber erschrecken, wenn man sie ernstlich fände. Die Philosophen wissen Grosses davon zu berichten; sie möchten gern vergessen machen, dass es einen unlösbaren Zusammenhang alles Geistigen giebt, ein feines, aber unzerreissbares Geflecht von Fäden und Beziehungen, das den Menschen umspinnt, solange er Mensch ist. Heraklit und Empedokles, und Alle ihresgleichen mit halbgöttlichen Ansprüchen und Erhabenheiten — sie waren Menschen, die sich von Menschen formten und färbten: in ihre weltfernste Einsamkeit trug irgend ein Wind noch einen Nachklang von Marktgeschwätz und Tagesgeräusch. Wieviel Heute und Gestern tönt in den ewigen Sprüchen aller Weisheit mit, wie Silber in broncenen Glocken; wieviel blöde Menschenmeinung, vom Philosophen in seine Höhe mitgenommen, bleibt dort über die nächste Jahrtausendwende liegen! Jeder Mensch verdankt jedem, und oft von seinem Besten; jeder Mensch befruchtet jeden, oft auf wunderlichen Zwischenwegen, durch Dritte und Vierte, die ohne Wissen und Willen den Gedankensamen von Geist zu Geist tragen. Es ist ergreifend und auf Reisen wohl zu beobachten, wie überall die Menschen von demselben Stoffe zehren, nämlich von einander, wie sich in jede Cultur die späteren theilen, sie aufsaugen und verwandeln, und von Morgen zu Abend hundert Ströme der Mittheilung laufen. In dieser Vielstimmigkeit aller Bildung und Gesittung eine misstönende Einzelheit bedeuten wollen, ist ebenso unnütz wie, umgekehrt, die Bereitwilligkeit, sich in die Harmonie einzufügen: keines von beiden steht in unserer Macht. Der Sonderling, der abseits bleibt, und der begeisterte Demagog, der sich zur „Hingabe an den Weltprocess“ und Vertretung des Gesammtwillens entschliesst — Beide nehmen ihren Platz unter Menschen ein, aber einen andern als sie denken.
\par

